\section{Robustness and Uncertainty}
\label{appendix:robustness}

In this section, we describe the details of the experimental setting and evaluation methods.

\subsection{Robustness}

We evaluate the model robustness to adversarial perturbations, occlusion and in-between samples using ImageNet trained models. For the base models, we use ResNet-50 structure and follow the settings in Section 4.1.1.
For comparison, we use ResNet-50 trained without any additional regularization or augmentation techniques, ResNet-50 trained by Mixup strategy, ResNet-50 trained by Cutout strategy and ResNet-50 trained by our proposed CutMix strategy.

\noindent\textbf{Fast Gradient Sign Method (FGSM): }
We employ Fast Gradient Sign Method (FGSM) \cite{fgsm} to generate adversarial samples.
For the given image $x$, the ground truth label $y$ and the noise size $\epsilon$, FGSM generates an adversarial sample as the following
\begin{equation}
    \hat x = x + \epsilon ~\texttt{sign}\left(\nabla_{x} L(\theta, x, y)\right),
\end{equation}
where $L(\theta, x, y)$ denotes a loss function, for example, cross entropy function.
In our experiments, we set the noise scale $\epsilon = 8/255$.

\noindent\textbf{Occlusion: }
For the given hole size $s$, we make a hole with width and height equals to $s$ in the center of the image.
For center occluded samples, we zeroed-out inside of the hole and for boundary occluded samples, we zeroed-out outside of the hole.
In our experiments, we test the top-1 ImageNet validation accuracy of the models with varying hole size from $0$ to $224$.

\noindent\textbf{In-between class samples: }
To generate in-between class samples, we first sample $50,000$ pairs of images from the ImageNet validation set. For generating Mixup samples, we generate a sample $x$ from the selected pair $x_A$ and $x_B$ by $x = \lambda x_A + (1-\lambda) x_B$. We report the top-1 accuracy on the Mixup samples by varying $\lambda$ from $0$ to $1$.
To generate CutMix in-between samples, we employ the center mask instead of the random mask. We follow the hole generation process used in the occlusion experiments. We evaluate the top-1 accuracy on the CutMix samples by varing hole size $s$ from $0$ to $224$.


\begin{table*}[]
\centering
\begin{tabular}{@{}lccc|ccc@{}}
\toprule
Method   & TNR at TPR 95\% & AUROC & Detection Acc.  & TNR at TPR 95\% & AUROC & Detection Acc. \\ \midrule
         & \multicolumn{3}{c|}{TinyImageNet} & \multicolumn{3}{c}{TinyImageNet (resize)}       \\ \midrule
Baseline & 43.0 (0.0)    & 88.9 (0.0)   & 81.3 (0.0) & 29.8 (0.0)    & 84.2 (0.0)   & 77.0 (0.0)  \\
Mixup    & 22.6 \redp{-20.4}  & 71.6 \redp{-17.3} & 69.8 \redp{-11.5} & 12.3 \redp{-17.5}  & 56.8 \redp{-27.4} & 61.0 \redp{-16.0} \\
Cutout   & 30.5 \redp{-12.5}  & 85.6 \redp{-3.3}  & 79.0 \redp{-2.3}  & 22.0 \redp{-7.8}   & 82.8 \redp{-1.4}  & 77.1 \greenp{+0.1}  \\
CutMix   & \textbf{57.1 \greenp{+14.1}}   & \textbf{92.4 \greenp{+3.5}}   & \textbf{85.0 \greenp{+3.7}}  & \textbf{55.4 \greenp{+25.6}}   & \textbf{91.9 \greenp{+7.7}}   & \textbf{84.5 \greenp{+7.5}}  \\ \midrule
         & \multicolumn{3}{c|}{LSUN (crop)} & \multicolumn{3}{c}{LSUN (resize)}\\ \midrule
Baseline & 34.6 (0.0)    & 86.5 (0.0)   & 79.5 (0.0) & 34.3 (0.0)    & 86.4 (0.0)   & 79.0 (0.0)  \\
Mixup    & 22.9 \redp{-11.7}  & 76.3 \redp{-10.2} & 72.3 \redp{-7.2} & 13.0 \redp{-21.3}  & 59.0 \redp{-27.4} & 61.8 \redp{-17.2} \\
Cutout   & 33.2 \redp{-1.4}   & 85.7 \redp{-0.8}  & 78.5 \redp{-1.0} & 23.7 \redp{-10.6}  & 84.0 \redp{-2.4}  & 78.4 \redp{-0.6} \\
CutMix   & \textbf{47.6 \greenp{+13.0}}   & \textbf{90.3 \greenp{+3.8}}   & \textbf{82.8 \greenp{+3.3}} & \textbf{62.8 \greenp{+28.5}}   & \textbf{93.7 \greenp{+7.3}}   & \textbf{86.7 \greenp{+7.7}}  \\ \midrule
    & &   \multicolumn{3}{c}{iSUN} \\ \midrule
Baseline & & 32.0 (0.0)    & \multicolumn{1}{c}{85.1 (0.0}   & 77.8 (0.0)  \\
Mixup   & & 11.8 \redp{-20.2}  & \multicolumn{1}{c}{57.0 \redp{-28.1}} & 61.0 \redp{-16.8}  \\
Cutout  & & 22.2 \redp{-9.8}   & \multicolumn{1}{c}{82.8 \redp{-2.3}}  & 76.8 \redp{-1.0} \\
CutMix  & & \textbf{60.1 \greenp{+28.1}}   & \multicolumn{1}{c}{\textbf{93.0 \greenp{+7.9}}}   & \textbf{85.7 \greenp{+7.9}} \\ \midrule
         & \multicolumn{3}{c|}{Uniform} & \multicolumn{3}{c}{Gaussian} \\ \midrule
Baseline & 0.0 (0.0)     & 89.2 (0.0)   & 89.2 (0.0) & 10.4 (0.0)    & 90.7 (0.0)   & 89.9 (0.0)   \\
Mixup    & 0.0 (0.0)     & 0.8 \redp{-88.4}  & 50.0 \redp{-39.2} & 0.0 \redp{-10.4}   & 23.4 \redp{-67.3} & 50.5 \redp{-39.4} \\
Cutout   & 0.0 (0.0)     & 35.6 \redp{-53.6} & 59.1 \redp{-30.1} & 0.0 \redp{-10.4}   & 24.3 \redp{-66.4} & 50.0 \redp{-39.9} \\
CutMix   & \textbf{100.0 \greenp{+100.0}} & \textbf{99.8 \greenp{+10.6}}  & \textbf{99.7 \greenp{+10.5}}  & \textbf{100.0 \greenp{+89.6}}  & \textbf{99.7 \greenp{+9.0}}   & \textbf{99.0 \greenp{+9.1}}  \\ \bottomrule
\end{tabular}
\caption{Out-of-distribution (OOD) detection results on TinyImageNet, LSUN, iSUN, Gaussian noise and Uniform noise using CIFAR-100 trained models. All numbers are in percents; higher is better.}
\label{table:ood-sup}
\end{table*}

\subsection{Uncertainty}

Deep neural networks are often overconfident in their predictions. For example, deep neural networks produce high confidence number even for random noise \cite{hendrycks2016baseline}. One standard benchmark to evaluate the overconfidence of the network is Out-of-distribution (OOD) detection proposed by \cite{hendrycks2016baseline}. The authors proposed a threshold-baed detector which solves the binary classification task by classifying in-distribution and out-of-distribution using the prediction of the given network. Recently, a number of reserchs are proposed to enhance the performance of the baseline detector \cite{liang2017odin, lee2018simple} but in this paper, we follow only the baseline detector algorithm without any input enhancement and temperature scaling \cite{liang2017odin}.

\noindent\textbf{Setup: }
We compare the OOD detector performance using CIFAR-100 trained models described in Section 4.1.2. For comparison, we use PyramidNet-200 model without any regularization method, PyramidNet-200 model with Mixup, PyramidNet-200 model with Cutout and PyramidNet-200 model with our proposed CutMix.

\noindent\textbf{Evaluation Metrics and Out-of-distributions: }
In this work, we follow the experimental setting used in \cite{hendrycks2016baseline, liang2017odin}.
To measure the performance of the OOD detector, we report the true negative rate (TNR) at 95\% true positive rate (TPR), the area under the receiver operating characteristic curve (AUROC) and detection accuracy of each OOD detector.
We use seven datasets for out-of-distribution: TinyImageNet (crop), TinyImageNet (resize), LSUN \cite{yu2015lsun} (crop), LSUN (resize), iSUN, Uniform noise and Gaussian noise.

\noindent\textbf{Results: } 
We report OOD detector performance to seven OODs in Table~\ref{table:ood-sup}. Overall, CutMix outperforms baseline, Mixup and Cutout. Moreover, we find that even though Mixup and Cutout outperform the classification performance, Mixup and Cutout largely degenerate the baseline detector performance.
Especially, for Uniform noise and Gaussian noise, Mixup and Cutout seriously impair the baseline performance while CutMix dramatically improves the performance.
From the experiments, we observe that our proposed CutMix enhances the OOD detector performance while Mixup and Cutout produce more overconfident predictions to OOD samples than the baseline.